% =====================================================================
% 04a -- Companion experimental program.
% =====================================================================
% Planning (% comments only):
% This section ties the review paper to the repo experiments. The core
% principle is: every critique should either be a theorem/literature
% claim or have a falsifiable probe with a toy control and a real-data
% test. The first real-data domain is VeriBench/Lean; MNIST is the first
% simple vision domain for order/global-validity effects.
% =====================================================================

\section{Companion Experimental Program}
\label{sec:experimental_program}

The previous sections separate the objections by layer and identify which ones are worth taking seriously.
This section describes the empirical program that keeps the review honest.
The goal is not to produce a single benchmark score in favor of or against autoregressive models.
The goal is to test each objection at the layer where it actually lives, first in a toy setting where the effect should be visible by construction, and then in a real domain where the effect may be masked, amplified, or reversed.

The coordination folder for this program is \texttt{experiments/03\_review\_paper/}.
The first concrete claim suite is \texttt{experiments/02\_ar\_pros\_cons/}.
The paper, blogs, and experiments are intended to move together: a claim becomes paper-ready only after it has a reproducible command, a saved result, a figure or table, and a caveat.

\subsection{Three experimental tiers}
\label{ssec:three_tiers}

\paragraph{Tier 1: toy controls.}
Toy controls are not substitutes for real data; they are positive controls.
For the LeCun compounding argument, a blind independent-error rollout should follow
\begin{equation}
  \PP[\text{success at length } T] = (1-\varepsilon)^T.
  \label{eq:lecun_toy_control}
\end{equation}
If a toy built to satisfy the assumptions of \cref{eq:lecun_toy_control} does not produce the geometric curve, the probe is broken.
The useful comparisons are then verifier-resampling and recoverable-error toys, where the same raw local error rate need not produce the same surviving sequence-level error rate.
The current toy harness lives in \texttt{experiments/02\_ar\_pros\_cons/toy/}.
We explicitly leave a review item for Elyas Obbad to propose a toy closer to \textsc{VeriBench}: a setting that shows both the claimed autoregressive failure mode and the real autoregressive advantages of exact likelihood and cheap local normalization.

\paragraph{Tier 2: VeriBench/Lean.}
VeriBench supplies the hard-checker domain this paper needs~\citep{daruru2024veribench,moura2021lean}.
The relevant outcome is not token accuracy but verifier survival: whether a generated Lean file, theorem, or proof passes the Lean kernel and associated checks.
This matters because the strongest form of the compounding argument assumes unrecoverability.
A verifier-backed system is designed precisely to reject, retry, and recover.

The current data splitter scans a local \texttt{\textasciitilde/veribench/veribench\_dataset} checkout and writes deterministic train/validation/test manifests at the task level.
The first real-data analysis should fit three models to pass probability as a function of length or proof-depth proxy:
\begin{enumerate}[leftmargin=2em]
  \item a geometric model, corresponding to \cref{eq:lecun_toy_control};
  \item a constant-pass model, corresponding to length-insensitive success;
  \item a recoverable Markov model, corresponding to nonzero probability of returning to the valid set after an error.
\end{enumerate}
If the recoverable model fits better, the conclusion is not that \cref{eq:lecun_toy_control} is algebraically false.
The conclusion is that it is the wrong model of a verifier-guided system.
If the geometric model fits even after retry/search/verifier feedback, then the compounding critique survives in exactly the domain where we expected it to fail.

\paragraph{Tier 3: a simple vision domain.}
The first vision domain can be \textsc{MNIST}~\citep{lecun1998gradient}.
The purpose is methodology rather than benchmark leadership.
Images can be treated as token sequences under raster, shuffled, patch, or learned orderings.
This makes it possible to compare:
\begin{itemize}
  \item raster-order autoregressive reconstruction;
  \item shuffled-order autoregressive reconstruction;
  \item masked or iterative refinement;
  \item an energy classifier or ranker over complete images.
\end{itemize}
The measurement is whether local prediction and ordering choices create measurable global-validity or classifier-agreement failures after controlling for parameter count and compute.
Later, MNIST should be replaced or supplemented by a richer vision dataset, but the first pass should remain small enough that methodological mistakes are visible.

\subsection{Integrated tests rather than isolated anecdotes}
\label{ssec:integrated_tests}

The isolated probes establish that mechanisms exist.
They do not establish that a mechanism moves the downstream metric once everything is trained together.
For that, the companion suite uses an integrated harness.
One run logs hidden-state rank, attention entropy, bottleneck rank deficit, invalid-token mass, spectral-norm product, and verifier pass rate through training.
A factorial ablation grid then toggles the commitments separately:
\begin{center}
\begin{tabular}{lll}
\toprule
Axis & Baseline & Alternative \\
\midrule
attention normalization & softmax & sigmoid/linear or energy-style attention \\
objective & MLE & margin/energy/verifier-aware objective \\
inference & blind AR rollout & verifier-guided retry/search \\
ordering & left-to-right & any-order/masked/iterative \\
\bottomrule
\end{tabular}
\end{center}
The scientific object is not just the main effect of each toggle but the interactions.
For example, replacing softmax may do little under MLE but matter under an energy or margin objective.
That is a different conclusion from ``softmax is unnecessary,'' and it is the kind of conclusion the review paper should be able to state cleanly.

\subsection{How results enter the paper}
\label{ssec:result_promotion}

A result should be promoted from experiment log to paper text only after it satisfies five conditions:
\begin{enumerate}[leftmargin=2em]
  \item a reproducible command exists in the experiment folder;
  \item JSON or tabular statistics are saved;
  \item a figure or table exists;
  \item uncertainty is reported with seeds or bootstrap intervals when applicable;
  \item the caveat is written before the interpretation.
\end{enumerate}
The blog series is allowed to be earlier and more exploratory.
The paper should be stricter.
This is why the first blog post can introduce the LeCun experiment before the \textsc{VeriBench} result is complete, but this section treats the result as pending until the real-data run lands in \texttt{FINDINGS.md}.
